% =====================================================================
\section{Fase de Producción}
\label{sec:produccion}

La fase de producción se desarrolló en las seis iteraciones planificadas en la
Tabla~\ref{tab:plan-iteraciones}. Siguiendo Mobile-D, cada iteración comenzó
con un \emph{día de planificación}, en el que se redactaron las historias de
usuario y sus criterios de aceptación; continuó con los \emph{días de
trabajo}, dedicados al diseño y la implementación; y cerró con un \emph{día de
liberación}, en el que se verificaron los criterios de aceptación y el
módulo se integró en la rama \archivo{develop}. Los resultados detallados de
la verificación se consolidan en la sección~\ref{sec:pruebas} y en el
Anexo~\ref{anx:casos}.

% ---------------------------------------------------------------------
\subsection{Iteración 1: Autenticación, onboarding y perfil}
\label{subsec:it1}

\subsubsection{Planificación}

\begin{table}[H]
\centering
\caption{Historias de usuario de la iteración 1}
\label{tab:hu-it1}
\begin{tabularx}{\textwidth}{L{1.2cm} Y L{1.5cm}}
\toprule
\textbf{ID} & \textbf{Historia de usuario} & \textbf{RF} \\
\midrule
HU-01 & Como usuario nuevo, quiero registrarme con mi nombre, correo y
contraseña para tener una cuenta propia. & RF1 \\
HU-02 & Como usuario, quiero iniciar sesión con mi cuenta de Google para no
recordar otra contraseña. & RF1 \\
HU-03 & Como usuario, quiero recuperar mi contraseña por correo si la
olvido. & RF1 \\
HU-04 & Como usuario recién registrado, quiero indicar mi objetivo y mi nivel
para recibir recomendaciones acordes. & RF2 \\
HU-05 & Como usuario, al terminar el registro quiero elegir entre generar una
rutina automática o armar la mía. & RF2 \\
HU-06 & Como usuario, quiero editar mis datos de perfil y cerrar sesión. &
RF12 \\
\bottomrule
\end{tabularx}
\end{table}

\subsubsection{Trabajo}

\paragraph{Controlador de autenticación.} Toda la interacción con Firebase
Authentication se concentró en \archivo{AuthController}. El registro crea la
cuenta y, a continuación, el documento de perfil \archivo{usuarios/\{uid\}}.
En el inicio de sesión con Google, Firebase crea la cuenta pero no el perfil,
por lo que el controlador invoca \archivo{asegurarPerfil}, que lo crea solo si
no existe (Código~\ref{lst:auth-registro}).

\begin{lstlisting}[language=TypeScript, caption={Registro e inicio de sesión con Google en \texttt{AuthController}}, label={lst:auth-registro}]
async registrar(datos: DatosRegistro): Promise<string> {
  const credencial = await createUserWithEmailAndPassword(
    auth, datos.correo, datos.contrasena);
  const uid = credencial.user.uid;
  await crearPerfilUsuario(uid, {
    nombre: datos.nombre, apellidos: datos.apellidos,
    genero: datos.genero, email: datos.correo,
  });
  return uid;
},

async iniciarSesionConGoogle(idToken: string): Promise<UsuarioSesion> {
  const credencial = GoogleAuthProvider.credential(idToken);
  const { user } = await signInWithCredential(auth, credencial);
  await AuthController.asegurarPerfil(user.uid, {
    nombre: user.displayName, email: user.email,
  });
  return aUsuarioSesion(user);
},
\end{lstlisting}

Los formularios se validan con esquemas Zod antes de llamar al controlador:
el registro exige nombre y apellido, un correo válido, una contraseña de al
menos seis caracteres y la coincidencia con su confirmación.

\paragraph{Estado de la sesión.} El estado global se gestiona con un
\emph{store} de Zustand persistido en el dispositivo, alimentado por dos
suscripciones montadas una sola vez en el diseño raíz: \archivo{useAuthListener},
que refleja los cambios de sesión de Firebase, y \archivo{usePerfilListener},
que observa en tiempo real el documento del perfil y calcula si está completo.
La regla de negocio es explícita: un perfil está completo cuando tiene
objetivo y nivel, los datos mínimos para generar una rutina. La
Figura~\ref{fig:secuencia-login} muestra cómo estos elementos colaboran en el
inicio de sesión.

\begin{figure}[H]
\centering
\begin{tikzpicture}[font=\scriptsize, x=2.45cm, y=0.62cm,
  part/.style={rectangle, draw=black!70, fill=green!8, minimum width=2.0cm,
    minimum height=0.6cm, align=center, font=\tiny\bfseries},
  msg/.style={-{Stealth[length=2mm]}, thick},
  ret/.style={-{Stealth[length=2mm]}, dashed, black!60}]
  \foreach \i/\n in {0/Usuario, 1/LoginScreen, 2/AuthController, 3/Firebase Auth, 4/Firestore, 5/Layout raíz} {
    \node[part] (p\i) at (\i,0) {\n};
    \draw[black!40, dashed] (\i,-0.5) -- (\i,-12.2);
  }
  \draw[msg] (0,-1.6) -- node[above]{correo y contraseña} (1,-1.6);
  \node[draw, fill=white, font=\tiny, anchor=west] at (1.05,-2.3) {valida con Zod};
  \draw[msg] (1,-3.4) -- node[above]{iniciarSesion()} (2,-3.4);
  \draw[msg] (2,-4.4) -- node[above]{signInWith...()} (3,-4.4);
  \draw[ret] (3,-5.4) -- node[above]{sesión (uid)} (2,-5.4);
  \draw[ret] (3,-6.6) -- node[above, pos=0.5]{onAuthStateChanged $\rightarrow$ store} (5,-6.6);
  \draw[msg] (5,-7.8) -- node[above]{observar usuarios/\{uid\}} (4,-7.8);
  \draw[ret] (4,-9.0) -- node[above]{perfil} (5,-9.0);
  \node[draw, fill=white, align=left, font=\scriptsize, anchor=north east] at (5.1,-9.6)
    {¿perfil completo?\\sí $\rightarrow$ \texttt{(app)}\\no $\rightarrow$ \texttt{(onboarding)}};
\end{tikzpicture}
\caption{Diagrama de secuencia del inicio de sesión}
\label{fig:secuencia-login}
\end{figure}

\paragraph{Onboarding.} El grupo \archivo{(onboarding)} tiene dos pantallas.
En \archivo{TargetSelectionScreen} el usuario elige uno de los cuatro
objetivos y uno de los tres niveles de experiencia, cada uno con una
descripción que lo ayuda a ubicarse (por ejemplo, ``Principiante: soy nuevo
en el gimnasio o llevo menos de 6 meses''). En \archivo{RoutineChoiceScreen}
decide si quiere generar una rutina automática, crear una personalizada o
continuar sin rutina. Al guardar el objetivo y el nivel en Firestore, el
listener del perfil cambia el estado y el diseño raíz habilita
\archivo{(app)}; la opción elegida queda registrada en el \emph{store} y la
pantalla de inicio la abre automáticamente.

\subsubsection{Liberación}

\begin{table}[H]
\centering
\caption{Criterios de aceptación de la iteración 1}
\label{tab:ca-it1}
\begin{tabularx}{\textwidth}{L{1.3cm} Y}
\toprule
\textbf{HU} & \textbf{Criterio de aceptación} \\
\midrule
HU-01 & Con datos válidos se crea la cuenta y el documento
\archivo{usuarios/\{uid\}}; con un correo inválido, contraseñas distintas o
de menos de seis caracteres se muestra el mensaje correspondiente. \\
HU-02 & El inicio con Google crea el perfil la primera vez y no lo
sobrescribe en accesos posteriores. \\
HU-03 & Se envía el correo de restablecimiento a la dirección ingresada. \\
HU-04 & Un usuario sin objetivo o sin nivel es dirigido al onboarding y no
puede acceder a \archivo{(app)}. \\
HU-05 & La opción elegida abre la pantalla de generación o de rutina
personalizada al llegar al inicio. \\
HU-06 & Los cambios del perfil se guardan; al cerrar sesión se vuelve a la
bienvenida. \\
\bottomrule
\end{tabularx}
\end{table}

\capturados{login}{Inicio de sesión}{registro}{Registro}
  {Pantallas de inicio de sesión y registro}{fig:cap-auth}
\capturados{onboarding-objetivo}{Objetivo y nivel}{onboarding-rutina}{Elección de rutina}
  {Pantallas del onboarding}{fig:cap-onboarding}

% ---------------------------------------------------------------------
\subsection{Iteración 2: Guía de máquinas y ejercicios}
\label{subsec:it2}

\subsubsection{Planificación}

\begin{table}[H]
\centering
\caption{Historias de usuario de la iteración 2}
\label{tab:hu-it2}
\begin{tabularx}{\textwidth}{L{1.2cm} Y L{1.5cm}}
\toprule
\textbf{ID} & \textbf{Historia de usuario} & \textbf{RF} \\
\midrule
HU-07 & Como usuario, quiero buscar una máquina por su nombre o por el músculo
que trabaja, aunque no escriba las tildes. & RF3 \\
HU-08 & Como usuario, quiero filtrar el catálogo por grupo muscular y por tipo
de equipo. & RF3 \\
HU-09 & Como principiante, quiero ver paso a paso cómo se usa una máquina, qué
errores evitar y cómo prevenir lesiones. & RF4 \\
HU-10 & Como usuario, quiero ver una animación de cada ejercicio con su
músculo objetivo y su dificultad. & RF5 \\
HU-11 & Como usuario, quiero ver cómo luce la máquina en distintos modelos
comerciales para reconocerla en mi gimnasio. & RF4 \\
\bottomrule
\end{tabularx}
\end{table}

\subsubsection{Trabajo}

\paragraph{Construcción del catálogo.} El contenido de la guía se preparó
fuera de la aplicación mediante un proceso de cuatro pasos
(Figura~\ref{fig:pipeline-catalogo}):

\begin{enumerate}
    \item \textbf{Selección de máquinas:} a partir del catálogo de un
    proveedor nacional se definieron 35 tipos de máquina, que agrupan 288
    modelos comerciales como variantes; se excluyeron los accesorios
    (agarres, barras y soportes).
    \item \textbf{Mapeo de ejercicios:} cada tipo se asoció manualmente con
    sus ejercicios en wger (98 ejercicios) y en free-exercise-db (91
    ejercicios). El mapeo se hizo a mano porque la coincidencia automática por
    nombre producía asociaciones incorrectas.
    \item \textbf{Contenido propio:} para cada máquina se redactaron la
    descripción, los músculos principales y secundarios, la guía de uso, los
    errores comunes y las pautas de prevención de lesiones, contrastando la
    información con guías de entrenamiento reconocidas sin copiar sus textos.
    \item \textbf{Animaciones:} se asociaron 95 animaciones GIF de WorkoutX,
    al menos una por máquina, descargadas mediante su API oficial y
    respetando sus condiciones de uso.
\end{enumerate}

\begin{figure}[H]
\centering
\begin{tikzpicture}[node distance=0.45cm, font=\footnotesize,
  p/.style={caja, text width=2.35cm, minimum height=1.4cm}]
  \node[p] (a) {Catálogo del proveedor\\(35 tipos)};
  \node[p, right=of a] (b) {Mapeo a wger y free-exercise-db};
  \node[p, right=of b] (c) {Fichas con contenido propio};
  \node[p, right=of c] (d) {GIF de WorkoutX (API)};
  \node[p, fill=orange!8, below=0.9cm of c] (e) {Script de carga (SDK Admin)};
  \node[p, fill=orange!8, left=1.2cm of e] (f) {Firestore:\\texto};
  \node[p, fill=orange!8, right=1.2cm of e] (g) {Storage:\\imágenes y GIF};
  \draw[flecha] (a)--(b); \draw[flecha] (b)--(c); \draw[flecha] (c)--(d);
  \draw[flecha] (d) |- (e); \draw[flecha] (e)--(f); \draw[flecha] (e)--(g);
\end{tikzpicture}
\caption{Proceso de construcción y carga del catálogo de máquinas}
\label{fig:pipeline-catalogo}
\end{figure}

El \emph{script} de carga es idempotente: puede ejecutarse varias veces sin
duplicar documentos. Guarda el texto en \archivo{maquinas/\{id\}} y los
archivos multimedia en Storage, registrando en el documento la ruta de cada
archivo. Cada ejercicio conserva los identificadores de su fuente de origen,
lo que da trazabilidad sobre qué contenido es propio y cuál proviene de
fuentes abiertas.

\paragraph{Repositorio con caché.} El catálogo es pequeño (35 documentos,
alrededor de 150~KB), por lo que \archivo{maquinaRepository} lo descarga una
sola vez por sesión, valida cada documento con un esquema Zod (los documentos
inválidos se omiten con un aviso, en lugar de romper la pantalla) y lo
mantiene en memoria. Las búsquedas y filtros se resuelven localmente, sin
nuevas consultas a la red (RNF5). El Código~\ref{lst:filtrar} muestra la regla
de filtrado del controlador, que normaliza el texto para ignorar mayúsculas y
tildes.

\begin{lstlisting}[language=TypeScript, caption={Regla de filtrado del catálogo en \texttt{MaquinasController}}, label={lst:filtrar}]
function normalizar(texto: string): string {
  return texto.normalize("NFD").replace(/[̀-ͯ]/g, "")
    .toLowerCase().trim();
}

filtrar(catalogo: Maquina[], filtros: FiltrosMaquinas): Maquina[] {
  const texto = filtros.busqueda ? normalizar(filtros.busqueda) : "";
  return catalogo.filter((m) => {
    if (filtros.categoria && m.categoria !== filtros.categoria) return false;
    if (filtros.tipo && m.tipo !== filtros.tipo) return false;
    if (!texto) return true;
    const campos = [m.nombre, m.categoria, ...m.musculosPrincipales,
      ...m.musculosSecundarios, ...m.ejercicios.map((e) => e.nombre),
      ...m.ejerciciosGif.map((e) => e.nombre)];
    return campos.some((campo) => normalizar(campo).includes(texto));
  });
},
\end{lstlisting}

\paragraph{Imágenes desde Storage.} Las rutas de Storage se convierten en URL
de descarga mediante \archivo{storageService}, que almacena en caché cada URL
ya resuelta. En la interfaz, el componente \archivo{ImagenStorage}, construido
sobre \archivo{expo-image}, muestra tanto imágenes estáticas como animaciones
GIF.

\paragraph{Pantallas.} La guía consta de dos pantallas:

\begin{itemize}
    \item \textbf{Catálogo} (\archivo{machines/index}): cuadrícula de dos
    columnas con buscador, fichas de filtro por grupo muscular y por tipo, y
    actualización por arrastre.
    \item \textbf{Ficha} (\archivo{machines/[id]}): portada, botón ``Ver en
    Realidad Aumentada'', descripción, músculos, ejercicios con su animación y
    nivel de dificultad, guía de uso numerada, errores comunes, sección de
    seguridad, galería de variantes y atribución de las fuentes.
\end{itemize}

\subsubsection{Liberación}

\begin{table}[H]
\centering
\caption{Criterios de aceptación de la iteración 2}
\label{tab:ca-it2}
\begin{tabularx}{\textwidth}{L{1.3cm} Y}
\toprule
\textbf{HU} & \textbf{Criterio de aceptación} \\
\midrule
HU-07 & Buscar ``biceps'' o ``prensa'' devuelve las mismas máquinas que
``bíceps'' o ``Prensa''. \\
HU-08 & Los filtros de grupo muscular y tipo se combinan entre sí y con la
búsqueda. \\
HU-09 & Cada una de las 35 fichas muestra guía de uso, errores comunes y
pautas de seguridad. \\
HU-10 & Cada máquina presenta al menos un ejercicio con su animación,
músculo objetivo y dificultad. \\
HU-11 & La galería muestra hasta ocho variantes comerciales de la máquina. \\
\bottomrule
\end{tabularx}
\end{table}

\capturados{maquinas-catalogo}{Catálogo con filtros}{maquinas-ficha}{Ficha de máquina}
  {Guía de máquinas}{fig:cap-maquinas}

% ---------------------------------------------------------------------
\subsection{Iteración 3: Escáner de máquinas con realidad aumentada}
\label{subsec:it3}

\subsubsection{Planificación}

\begin{table}[H]
\centering
\caption{Historias de usuario de la iteración 3}
\label{tab:hu-it3}
\begin{tabularx}{\textwidth}{L{1.2cm} Y L{1.5cm}}
\toprule
\textbf{ID} & \textbf{Historia de usuario} & \textbf{RF} \\
\midrule
HU-12 & Como usuario, quiero enfocar una máquina con la cámara y ver su
información sobre la imagen real. & RF6 \\
HU-13 & Como usuario, si la máquina no se reconoce, quiero elegirla de una
lista sin salir del escáner. & RF6 \\
HU-14 & Como usuario, desde la ficha de una máquina quiero verla en modo de
realidad aumentada. & RF6 \\
\bottomrule
\end{tabularx}
\end{table}

\subsubsection{Trabajo}

\paragraph{Diseño.} El escáner (\archivo{machines/escaner}) se modeló como una
máquina de estados de cinco fases (Figura~\ref{fig:estados-escaner}). La
cámara se obtiene con \archivo{expo-camera}, solicitando el permiso con un
mensaje que explica su propósito. Durante el escaneo se dibuja un marco de
encuadre sobre la imagen en vivo; al confirmar la máquina, su nombre, tipo,
músculos y ejercicios con su dificultad se superponen sobre la cámara, junto
con un acceso a la ficha completa.

\begin{figure}[H]
\centering
\begin{tikzpicture}[node distance=2.0cm and 1.8cm, font=\footnotesize,
  est/.style={rectangle, rounded corners=8pt, draw=black!70, fill=green!8,
    minimum width=2.3cm, minimum height=0.8cm, align=center}]
  \node[est] (ini) at (0,0) {inicio};
  \node[est] (esc) at (3.8,0) {escaneando};
  \node[est] (ana) at (7.6,0) {analizando};
  \node[est] (eli) at (7.6,-2.6) {eligiendo};
  \node[est, fill=green!25] (det) at (3.8,-2.6) {detectada};
  \draw[flecha] (ini) -- node[above, font=\scriptsize]{permiso} (esc);
  \draw[flecha] (esc) -- node[above, font=\scriptsize]{capturar} (ana);
  \draw[flecha] (ana) -- node[right, font=\scriptsize, align=left]{sin modelo o\\confianza $<0.7$} (eli);
  \draw[flecha] (ana) -- node[above, sloped, font=\scriptsize, pos=0.45]{confianza $\geq 0.7$} (det);
  \draw[flecha] (eli) -- node[below, font=\scriptsize]{el usuario confirma} (det);
  \draw[flecha] (det.north) -- node[left, font=\scriptsize]{escanear otra} (esc.south);
  \draw[flecha] (ini.south) |- node[pos=0.75, below, font=\scriptsize]{desde la ficha (\texttt{?maquinaId})} (det.west);
\end{tikzpicture}
\caption{Diagrama de estados del escáner de máquinas}
\label{fig:estados-escaner}
\end{figure}

\paragraph{Punto de extensión para el reconocimiento automático.} La
identificación automática de la máquina se aisló en una única función,
\archivo{reconocerMaquina}, en la capa de servicios
(Código~\ref{lst:reconocimiento}). Su contrato recibe la fotografía y devuelve
el identificador de la máquina del catálogo con un valor de confianza entre 0
y 1. El controlador solo acepta la predicción si supera el umbral de 0.7; en
caso contrario, o si no hay modelo, la vista pasa al estado \emph{eligiendo} y
el usuario confirma la máquina en una lista con búsqueda.

\begin{lstlisting}[language=TypeScript, caption={Contrato del reconocimiento de máquinas y su uso en el controlador}, label={lst:reconocimiento}]
// services/ai/reconocimientoMaquinas.ts
export interface ResultadoReconocimiento {
  maquinaId: string;
  confianza: number; // 0..1
}
export const CONFIANZA_MINIMA = 0.7;
export async function reconocerMaquina(
  _fotoUri: string,
): Promise<ResultadoReconocimiento | null> {
  return null; // sin modelo entrenado: la Vista pide confirmación
}

// controllers/MaquinasController.ts
async reconocerDesdeFoto(fotoUri: string): Promise<Maquina | null> {
  const resultado = await reconocerMaquina(fotoUri);
  if (!resultado || resultado.confianza < CONFIANZA_MINIMA) return null;
  return obtenerMaquina(resultado.maquinaId);
},
\end{lstlisting}

Con este diseño, incorporar un clasificador de imágenes con TensorFlow Lite
consiste en implementar esa sola función; ni el controlador ni la vista deben
modificarse. Durante la iteración se constató que las animaciones de WorkoutX
no pueden usarse para entrenar modelos sin permiso escrito de su proveedor,
por lo que el conjunto de datos para el clasificador deberá construirse con
fotografías propias o con autorización (sección~\ref{sec:recomendaciones}).

\subsubsection{Liberación}

\begin{table}[H]
\centering
\caption{Criterios de aceptación de la iteración 3}
\label{tab:ca-it3}
\begin{tabularx}{\textwidth}{L{1.3cm} Y}
\toprule
\textbf{HU} & \textbf{Criterio de aceptación} \\
\midrule
HU-12 & Con el permiso concedido se muestra la cámara en vivo con el marco
de encuadre, y la ficha resumida se superpone sobre la imagen. \\
HU-13 & Sin reconocimiento automático, la lista con búsqueda permite elegir
cualquiera de las 35 máquinas y pasar al estado \emph{detectada}. \\
HU-14 & Al abrir el escáner desde una ficha, la máquina se muestra
directamente sobre la cámara. \\
\bottomrule
\end{tabularx}
\end{table}

\capturados{escaner-camara}{Escaneo}{escaner-detectada}{Ficha superpuesta}
  {Escáner de máquinas con realidad aumentada}{fig:cap-escaner}

% ---------------------------------------------------------------------
\subsection{Iteración 4: Rutinas personalizadas y planificación}
\label{subsec:it4}

\subsubsection{Planificación}

\begin{table}[H]
\centering
\caption{Historias de usuario de la iteración 4}
\label{tab:hu-it4}
\begin{tabularx}{\textwidth}{L{1.2cm} Y L{1.5cm}}
\toprule
\textbf{ID} & \textbf{Historia de usuario} & \textbf{RF} \\
\midrule
HU-15 & Como usuario, quiero que la aplicación genere una rutina semanal
adecuada a mi objetivo y mi nivel. & RF7 \\
HU-16 & Como usuario avanzado, quiero armar mi propia rutina día por día,
eligiendo los ejercicios y sus series, repeticiones y descansos. & RF8 \\
HU-17 & Como usuario, quiero agregar a mi rutina un ejercicio que veo en la
ficha de una máquina. & RF8 \\
HU-18 & Como usuario, quiero tener varias rutinas guardadas, elegir cuál es
la activa y eliminar las que ya no uso. & RF8 \\
HU-19 & Como usuario, al tocar un ejercicio de mi rutina quiero abrir la
ficha de su máquina. & RF8 \\
\bottomrule
\end{tabularx}
\end{table}

\subsubsection{Trabajo}

\paragraph{Generador de rutinas basado en reglas.} El componente inteligente
del sistema es un generador de rutinas basado en conocimiento, cuya estructura
sigue la de un sistema experto (sección~\ref{sec:mt-ia}):

\begin{itemize}
    \item La \textbf{memoria de trabajo} contiene los hechos del usuario: su
    objetivo y su nivel, tomados del perfil.
    \item La \textbf{base de conocimiento} está formada por tres conjuntos de
    reglas derivadas de las pautas de prescripción del ejercicio
    (Tablas~\ref{tab:reglas-frecuencia}, \ref{tab:reglas-parametros}
    y~\ref{tab:reglas-division}).
    \item El \textbf{motor de inferencia} aplica las reglas en secuencia y
    produce un plan semanal, que luego se completa con ejercicios reales de la
    API de wger.
\end{itemize}

La primera regla determina la frecuencia semanal según la experiencia,
siguiendo la progresión recomendada por \citet{kraemer2004}.

\begin{table}[H]
\centering
\caption{Reglas de frecuencia semanal según el nivel}
\label{tab:reglas-frecuencia}
\begin{tabular}{L{6.5cm} C{3cm}}
\toprule
\textbf{Regla} & \textbf{Días por semana} \\
\midrule
SI nivel = principiante & 3 \\
SI nivel = intermedio & 4 \\
SI nivel = avanzado & 5 \\
\bottomrule
\end{tabular}
\fuente{Archivo \archivo{utils/generadorRutinas.ts}, con base en \citet{kraemer2004}.}
\end{table}

La segunda regla fija las variables de carga según el objetivo, dentro de los
rangos de la Tabla~\ref{tab:variables-objetivo}: más series, menos
repeticiones y descansos largos para la hipertrofia; menos carga, más
repeticiones y descansos cortos para la resistencia y la pérdida de peso.

\begin{table}[H]
\centering
\caption{Reglas de variables de carga según el objetivo}
\label{tab:reglas-parametros}
\begin{tabular}{L{5.4cm} C{1.6cm} C{2.7cm} C{2.1cm}}
\toprule
\textbf{Regla} & \textbf{Series} & \textbf{Repeticiones} & \textbf{Descanso} \\
\midrule
SI objetivo = ganar masa muscular & 4 & 9 & 90 s \\
SI objetivo = mantenimiento & 3 & 12 & 60 s \\
SI objetivo = perder peso & 3 & 15 & 45 s \\
SI objetivo = resistencia & 3 & 18 & 30 s \\
\bottomrule
\end{tabular}
\fuente{Archivo \archivo{utils/generadorRutinas.ts}, con base en
\citet{kraemer2004} y \citet{schoenfeld2010}.}
\end{table}

La tercera regla define la división semanal combinando la frecuencia con el
objetivo. Se utilizan cinco bloques de día: \emph{empuje} (pecho, hombros y
brazos), \emph{tirón} (espalda y abdominales), \emph{pierna} (piernas y
pantorrillas), \emph{cardio} (cardio y abdominales) y \emph{cuerpo completo}.
Cuando el objetivo es perder peso o mejorar la resistencia, la regla
introduce un día de cardio en lugar de un día de fuerza.

\begin{table}[H]
\centering
\caption{Reglas de división semanal}
\label{tab:reglas-division}
\begin{tabularx}{\textwidth}{C{1.3cm} Y Y}
\toprule
\textbf{Días} & \textbf{Objetivo de fuerza (ganar masa, mantenimiento)} &
\textbf{Objetivo con énfasis cardio (perder peso, resistencia)} \\
\midrule
3 & Empuje -- Tirón -- Pierna & Empuje -- Cardio -- Pierna \\
4 & Empuje -- Tirón -- Pierna -- Cuerpo completo & Empuje -- Tirón -- Cardio
-- Pierna \\
5 & Empuje -- Tirón -- Pierna -- Empuje -- Tirón & Empuje -- Tirón -- Pierna
-- Cardio -- Cuerpo completo \\
\bottomrule
\end{tabularx}
\fuente{Archivo \archivo{utils/generadorRutinas.ts}.}
\end{table}

Cada bloque indica de qué categorías de ejercicios de wger tomar los
ejercicios y cuántos (entre cuatro y seis). El controlador reparte esa
cantidad entre las categorías del bloque lo más parejo posible, consulta los
ejercicios de cada categoría y elige al azar, sin repetir, los necesarios,
lo que da variedad entre rutinas generadas con los mismos datos. La
Figura~\ref{fig:flujo-generador} resume el proceso completo.

\begin{figure}[H]
\centering
\begin{tikzpicture}[node distance=0.42cm, font=\footnotesize,
  p/.style={caja, text width=4.6cm}]
  \node[p, fill=green!8] (a) {Hechos: objetivo y nivel del perfil};
  \node[p, below=of a] (b) {Validar con Zod (\archivo{generarRutinaIASchema})};
  \node[p, below=of b] (c) {Regla 1: días por semana según nivel};
  \node[p, below=of c] (d) {Regla 2: series, repeticiones y descanso según objetivo};
  \node[p, below=of d] (e) {Regla 3: bloques de la semana según días y objetivo};
  \node[p, below=of e] (f) {Por cada día: repartir cantidad entre categorías y
    elegir ejercicios de wger al azar};
  \node[p, fill=orange!8, below=of f] (g) {Guardar en \archivo{usuarios/\{uid\}/rutinas} y marcar como activa};
  \foreach \x/\y in {a/b,b/c,c/d,d/e,e/f,f/g} \draw[flecha] (\x)--(\y);
  \node[draw, dashed, rounded corners, fit=(c)(d)(e), inner sep=5pt,
    label={[font=\scriptsize\itshape]right:base de conocimiento}] {};
\end{tikzpicture}
\caption{Flujo del generador de rutinas basado en reglas}
\label{fig:flujo-generador}
\end{figure}

El Código~\ref{lst:generador} muestra el punto de entrada del motor de reglas.
Es una función pura, sin acceso a la red ni a la base de datos, lo que la hace
fácil de verificar y de reemplazar en el futuro por un modelo aprendido sin
modificar el controlador. El archivo completo se incluye en el
Anexo~\ref{anx:codigo}.

\begin{lstlisting}[language=TypeScript, caption={Punto de entrada del generador basado en reglas}, label={lst:generador}]
const PARAMETROS_POR_OBJETIVO: Record<ObjetivoUsuario, ParametrosSerie> = {
  perder_peso:   { series: 3, repeticiones: 15, descansoSegundos: 45 },
  ganar_masa:    { series: 4, repeticiones: 9,  descansoSegundos: 90 },
  mantenimiento: { series: 3, repeticiones: 12, descansoSegundos: 60 },
  resistencia:   { series: 3, repeticiones: 18, descansoSegundos: 30 },
};

const DIAS_POR_NIVEL: Record<NivelExperiencia, number> = {
  principiante: 3, intermedio: 4, avanzado: 5,
};

export function generarPlanHeuristico(
  objetivo: ObjetivoUsuario, nivel: NivelExperiencia,
): PlanHeuristico {
  const diasPorSemana = DIAS_POR_NIVEL[nivel];
  return {
    diasPorSemana,
    parametrosSerie: PARAMETROS_POR_OBJETIVO[objetivo],
    bloques: bloquesPorDias(diasPorSemana, objetivo),
  };
}
\end{lstlisting}

\paragraph{Rutina personalizada y planificación.} Para el objetivo OE4 se
construyó el planificador \archivo{routines/personalizada}, donde el usuario
agrega días, asigna a cada uno un enfoque y elige ejercicios desde un modal
de búsqueda, ajustando series, repeticiones y descanso. Los datos se validan
con Zod con límites razonables (de 1 a 7 días, de 1 a 20 series, de 1 a 100
repeticiones y hasta 10 minutos de descanso). Desde la ficha de cualquier
máquina, el botón ``Agregar a Mi Rutina'' permite añadir uno de sus
ejercicios a un día de la rutina activa; el controlador rechaza duplicados en
el mismo día.

La pantalla \archivo{routines/index} presenta el plan semanal de la rutina
activa como una lista de días expandibles y permite cambiar la rutina activa,
eliminar rutinas y crear nuevas. Las rutinas se observan en tiempo real con
\archivo{onSnapshot}, por lo que cualquier cambio se refleja sin recargar
(RNF4).

\subsubsection{Liberación}

\begin{table}[H]
\centering
\caption{Criterios de aceptación de la iteración 4}
\label{tab:ca-it4}
\begin{tabularx}{\textwidth}{L{1.3cm} Y}
\toprule
\textbf{HU} & \textbf{Criterio de aceptación} \\
\midrule
HU-15 & Para cada una de las doce combinaciones de objetivo y nivel, la
rutina generada tiene el número de días, la división y las variables de
carga de las Tablas~\ref{tab:reglas-frecuencia}
a~\ref{tab:reglas-division}. \\
HU-16 & Se puede crear una rutina de 1 a 7 días; un día sin ejercicios o
valores fuera de rango impiden guardar y muestran el motivo. \\
HU-17 & El ejercicio se agrega al día elegido; si ya existe en ese día, se
informa y no se duplica. \\
HU-18 & La rutina marcada como activa es la que se muestra en el inicio. \\
HU-19 & Al tocar un ejercicio asociado a una máquina se abre su ficha. \\
\bottomrule
\end{tabularx}
\end{table}

\capturados{rutina-generada}{Rutina generada}{rutina-personalizada}{Planificador personalizado}
  {Módulo de rutinas}{fig:cap-rutinas}

% ---------------------------------------------------------------------
\subsection{Iteración 5: Nutrición y calculadora de IMC}
\label{subsec:it5}

\subsubsection{Planificación}

\begin{table}[H]
\centering
\caption{Historias de usuario de la iteración 5}
\label{tab:hu-it5}
\begin{tabularx}{\textwidth}{L{1.2cm} Y L{1.5cm}}
\toprule
\textbf{ID} & \textbf{Historia de usuario} & \textbf{RF} \\
\midrule
HU-20 & Como usuario, quiero ingresar mi peso y mi estatura y conocer mi IMC
y su clasificación. & RF9 \\
HU-21 & Como usuario, quiero que cada cálculo quede guardado en mi
historial. & RF9 \\
\bottomrule
\end{tabularx}
\end{table}

\subsubsection{Trabajo}

De acuerdo con el alcance, el control calórico se resolvió mediante una
calculadora de IMC. El módulo fue el primero en implementarse respetando las
cuatro capas y sirvió como referencia para el resto: la regla de cálculo es
una función pura en \archivo{utils/imcCalculator.ts}
(Código~\ref{lst:imc}), que implementa la Ecuación~\ref{eq:imc} y la
clasificación de la Tabla~\ref{tab:imc-oms}; el controlador
\archivo{NutritionController} valida los datos (peso entre 10 y 500~kg y
estatura entre 50 y 300~cm), calcula, clasifica y guarda el registro en
\archivo{usuarios/\{uid\}/registrosNutricionales}; y la vista solo muestra el
resultado con una escala de colores de las cuatro categorías.

\begin{lstlisting}[language=TypeScript, caption={Cálculo y clasificación del IMC}, label={lst:imc}]
export function calcularIMC(pesoKg: number, alturaCm: number): number {
  const alturaM = alturaCm / 100;
  return Number((pesoKg / (alturaM * alturaM)).toFixed(1));
}

export function clasificarIMC(imc: number): BmiCategory {
  if (imc < 18.5) return "BAJO";
  if (imc < 25) return "NORMALIDAD";
  if (imc < 30) return "SOBREPESO";
  return "OBESIDAD";
}
\end{lstlisting}

Una decisión de diseño fue que un fallo al guardar no impida mostrar el
resultado: el cálculo no depende de la red, de modo que el controlador
registra el error y devuelve igualmente el IMC al usuario. El modelo
\archivo{RegistroNutricional} y su esquema de validación contemplan también
el registro de ingesta (alimento, calorías y cantidad), con su regla de
seguridad correspondiente, lo que deja preparada la ampliación a un diario de
alimentación.

\subsubsection{Liberación}

\begin{table}[H]
\centering
\caption{Criterios de aceptación de la iteración 5}
\label{tab:ca-it5}
\begin{tabularx}{\textwidth}{L{1.3cm} Y}
\toprule
\textbf{HU} & \textbf{Criterio de aceptación} \\
\midrule
HU-20 & Con 70~kg y 175~cm se obtiene un IMC de 22.9, categoría normal;
valores fuera de rango muestran un mensaje de validación. \\
HU-21 & Cada cálculo crea un documento en
\archivo{registrosNutricionales} con la fecha del servidor. \\
\bottomrule
\end{tabularx}
\end{table}

\captura{nutricion-imc}{Calculadora de IMC}{fig:cap-imc}

% ---------------------------------------------------------------------
\subsection{Iteración 6: Progreso físico y reportes}
\label{subsec:it6}

\subsubsection{Planificación}

\begin{table}[H]
\centering
\caption{Historias de usuario de la iteración 6}
\label{tab:hu-it6}
\begin{tabularx}{\textwidth}{L{1.2cm} Y L{1.5cm}}
\toprule
\textbf{ID} & \textbf{Historia de usuario} & \textbf{RF} \\
\midrule
HU-22 & Como usuario, quiero registrar mi peso periódicamente sin volver a
escribir mi estatura cada vez. & RF10 \\
HU-23 & Como usuario, quiero ver un gráfico de la evolución de mi peso. &
RF11 \\
HU-24 & Como usuario, quiero saber cuánto ha variado mi peso desde el primer
registro y cuál es mi categoría de IMC actual. & RF11 \\
HU-25 & Como usuario, quiero ver en el inicio un resumen de mi rutina activa
y de mis registros. & RF11 \\
\bottomrule
\end{tabularx}
\end{table}

\subsubsection{Trabajo}

\paragraph{Registro del progreso.} \archivo{ProgresoController} valida el
peso (entre 10 y 500~kg), la estatura opcional y unas notas de hasta 500
caracteres; si hay estatura, calcula el IMC reutilizando
\archivo{imcCalculator}, y guarda el registro en
\archivo{usuarios/\{uid\}/progresoFisico}. Para evitar que el usuario
reescriba su estatura, la pantalla la precarga desde el perfil.

\paragraph{Reporte.} La pantalla de progreso se suscribe en tiempo real a los
registros y deriva de ellos el reporte, sin almacenar datos calculados: un
gráfico de líneas con los últimos ocho registros, construido con
\archivo{react-native-chart-kit}; la variación del peso entre el primer y el
último registro; la categoría de IMC del último registro; y el historial de
mediciones. No se creó una colección de reportes porque el
reporte es una vista derivada del historial: calcularlo al vuelo garantiza que
siempre refleje los datos actuales y evita duplicarlos.

\paragraph{Inicio.} La pantalla de inicio reúne, en tiempo real, el saludo
con el nombre del usuario, la rutina activa con su número de ejercicios, el
número de registros de progreso y accesos directos al escáner, a la guía de
máquinas y a las rutinas.

\subsubsection{Liberación}

\begin{table}[H]
\centering
\caption{Criterios de aceptación de la iteración 6}
\label{tab:ca-it6}
\begin{tabularx}{\textwidth}{L{1.3cm} Y}
\toprule
\textbf{HU} & \textbf{Criterio de aceptación} \\
\midrule
HU-22 & La estatura del perfil aparece precargada; al guardar se crea el
registro con su IMC. \\
HU-23 & El gráfico se actualiza inmediatamente al agregar un registro y
muestra como máximo los ocho últimos. \\
HU-24 & Con registros de 80~kg y 77.5~kg se muestra una variación de
$-2.5$~kg; la categoría corresponde al IMC del último registro. \\
HU-25 & El inicio muestra la rutina activa y el número real de registros. \\
\bottomrule
\end{tabularx}
\end{table}

\capturados{progreso-reporte}{Reporte de progreso}{inicio}{Pantalla de inicio}
  {Progreso físico y pantalla de inicio}{fig:cap-progreso}
